\mmedAbsLabel{背景：}
跨媒介平台的快速更替正在重塑东亚娱乐产业的职业晋升路径。传统杂志、电视剧、综艺节目与视频平台在曝光机制、互动深度和风险传播速度上存在显著差异，使艺人的职业转型不再只是“工作类型变化”，而成为“平台规则、公众期待与个人品牌韧性”的联合调整过程。以佐佐木希为代表的长期活跃艺人提供了观察这一过程的公开案例。

\mmedAbsLabel{目的：}
构建一套适合公开案例研究的艺人职业转型分析框架，总结跨媒介转型的阶段特征、关键触发因素与形象韧性机制，并验证常规文本编码、事件时间线和简单量化指标在示例模板中的排版适配性。

\mmedAbsLabel{方法：}
基于公开新闻报道、访谈资料、节目单与社交媒体时间戳，建立“阶段划分-事件编码-韧性评分”三层分析框架。首先按出道、扩张、危机与重启四个阶段整理关键事件；随后将平台迁移、角色变化、公众反馈与自主表达纳入统一编码表；最后使用描述统计、对照表与示意图展示案例在不同阶段的媒介能力变化。

\mmedAbsLabel{结果：}
示例分析显示：第一，职业转型并非线性升级，而是以“平台机会窗口”驱动的阶段性重组；第二，传统媒体积累的可见度能够为后续影视与综艺转型提供起跳资源，但真正决定长期稳定性的，是危机发生后能否在新媒体中重建可控叙事；第三，形象韧性可以拆解为沉默期、修复期与重启期三个连续阶段，各阶段对应不同的话语策略与平台组合。基于公开案例构建的表格、长表格、图表和缩略词表均可在当前模板中稳定渲染。

\mmedAbsLabel{结论：}
跨媒介职业转型应被理解为“资源再配置+叙事再组织+平台再选择”的综合过程。对硕士论文写作而言，公开案例足以支撑完整的结构化示范，而无需依赖任何真实病历、真实受试者信息或未公开正文内容。

\mmedAbsLabel{关键词：}
跨媒介转型；案例研究；形象韧性；东亚娱乐产业；公开示例模板
